%% Uji gaya sitasi: chicago — Chicago (hal. 92 butir 3d)
%% Kompilasi: TEXINPUTS=../../:../../bst: BSTINPUTS=../../bst:
%%            pdflatex chicago && bibtex chicago && pdflatex chicago (2x)
%% (direktori bst/ harus terjangkau TEXINPUTS untuk pemeriksaan berkas gaya oleh kelas,
%%  dan BSTINPUTS supaya bibtex menemukan unnes-apalike-apa.bst)
\documentclass[sitasi=chicago,logodir=../../logo,logo=logo-placeholder]{unnes-ta}

\identitas{
  judul    = {Uji Gaya Sitasi chicago},
  nama     = {Nama Mahasiswa},
  nim      = {1234567890},
  prodi    = {Pendidikan Teknik Informatika},
  fakultas = {Fakultas Teknik},
  tahun    = {2026},
  pembimbinga = {Dr. Pembimbing Satu, M.Kom.},
}

\begin{document}
\bagianutama

\chapter{CONTOH SITASI}
\section{Bentuk Sitasi dalam Teks}
Penelitian campuran dibahas oleh \citet{creswell2018research} secara mendalam.
Kajian sebelumnya menegaskan hal yang sama \citep{martin2014write}.
Pendapat serupa disampaikan pada bagian tertentu \citep[955]{martin2014write}.
Rujukan bernomor mengikuti \cite{geraldi2017project}.
Dua sumber lain juga sejalan \citep{febriani2020memahami,shils1993etika}.
Bentuk MLA diuji tersendiri \citemla[23]{patten2017understanding}.
Buku metodologi seni dipakai sebagai acuan \citep{rohidi2012research}.
Regulasi akademik menjadi dasar \citep{permendikbudristek2023} dan
\citep{unnes2024panduanakademik}.

\daftarpustaka{../../referensi/contoh}

\end{document}
